\section{USAMO 2009}
        \begin{enumerate}
        	\item (\textit{USAMO 2009}) Given circles $\omega_1$ and $\omega_2$ intersecting at points $X$ and $Y$, let $\ell_1$ be a line through the center of $\omega_1$ intersecting $\omega_2$ at points $P$ and $Q$ and let $\ell_2$ be a line through the center of $\omega_2$ intersecting $\omega_1$ at points $R$ and $S$.  Prove that if $P, Q, R$ and $S$ lie on a circle then the center of this circle lies on line $XY$.


 



\textbf{Solution 1}

Let $\omega_3$ be the \href{https://artofproblemsolving.com/wiki/index.php?title=Circumcircle}{circumcircle} of $PQRS$, $r_i$ to be the radius of $\omega_i$, and $O_i$ to be the center of the circle $\omega_i$, where $i \in \{1,2,3\}$. Note that $SR$ and $PQ$ are the \href{https://artofproblemsolving.com/wiki/index.php?title=Radical\_axis}{radical axises} of $O_1$ , $O_3$ and $O_2$ , $O_3$ respectively. Hence, by \href{https://artofproblemsolving.com/wiki/index.php?title=Power\_of\_a\_point}{power of a point}(the power of $O_1$ can be expressed using circle $\omega_2$ and $\omega_3$ and the power of $O_2$ can be expressed using circle $\omega_1$ and $\omega_3$),
 \[O_1O_2^2 - r_2^2 = O_1O_3^2 - r_3^2\]
 \[O_2O_1^2 - r_1^2 = O_2O_3^2 - r_3^2\]Subtracting these two equations yields that $O_1O_3^2 - r_1^2 = O_2O_3^2 - r_2^2$, so $O_3$ must lie on the \href{https://artofproblemsolving.com/wiki/index.php?title=Radical\_axis}{radical axis} of $\omega_1$ , $\omega_2$.

 One needs to be careful of when $O_1$, $O_2$, and $O_3$ lie on a line. 

 


	\item (\textit{USAMO 2009}) Let $n$ be a positive integer.  Determine the size of the largest subset of $\{ - n, - n + 1, \ldots , n - 1, n\}$ which does not contain three elements $a, b, c$ (not necessarily distinct) satisfying $a + b + c = 0$.


 



\textbf{Solution 1}

Let $S$ be a subset of $\{-n,-n+1,\dots,n-1,n\}$ of largest size satisfying $a+b+c\neq 0$ for all $a,b,c\in S$. First, observe that $0\notin S$. Next note that $|S|\geq \lceil n/2\rceil$, by observing that the set of all the odd numbers in $\{-n,-n+1,\dots,n-1,n\}$ works. To prove that $|S|\leq \lceil n/2\rceil$, it suffices to only consider even $n$, because the statement for $2k$ implies the statement for $2k-1$ as well. So from here forth, assume $n$ is even.

 For any two sets $A$ and $B$, denote by $A+B$ the set $\{a+b\mid a\in A,b\in B\}$, and by $-A$ the set $\{-a\mid a\in A\}$. Also, let $A_+$ denote $A\cap\{1,2,\dots\}$ and $A_-$ denote $A\cap\{-1,-2,\dots\}$. First, we present a lemma:

 \textbf{Lemma 1}: Let $A$ and $B$ be two sets of integers. Then $|A+B|\geq|A|+|B|-1$.

 \textit{Proof}: Write $A=\{a_1,\dots,a_n\}$ and $B=\{b_1,\dots,b_m\}$ where $a_1<\dots<a_n$ and $b_1<\dots<b_m$. Then $a_1+b_1,a_1+b_2,\dots,a_1+b_m,a_2+b_m,\dots,a_n+b_m$ is a strictly increasing sequence of $n+m-1$ integers in $A+B$.

 Now, we consider two cases:

 \textbf{Case 1}: One of $n,-n$ is not in $S$. Without loss of generality, suppose $-n\notin S$. Let $U=\{-n+1,-n+2,\dots,n-2,n-1,n\}$ (a set with $2n$ elements), so that $S\subseteq U$ by our assumption. Now, the condition that $a+b+c\neq 0$ for all $a,b,c\in S$ implies that $-S\cap(S_++S_-)=\emptyset$. Since any element of $S_++S_-$ has absolute value at most $n-1$, we have $S_++S_-\subseteq \{-n+1,-n+2,\dots,n-2,n-1\}\subseteq U$. It follows that $S_++S_-\subseteq U\setminus -S$, so $|S_++S_-|\leq 2n-|S|$. However, by Lemma 1, we also have $|S_++S_-|\geq |S_+|+|S_-|-1=|S|-1$.  Therefore, we must have $|S|-1\leq 2n-|S|$, or $2|S|\leq 2n+1$, or $|S|\leq n$.

 \textbf{Case 2}: Both $n$ and $-n$ are in $S$. Then $n/2$ and $-n/2$ are not in $S$, and at most one of each of the pairs $\{1,n-1\},\{2,n-2\},\ldots,\{\frac n2-1,\frac n2+1\}$ and their negatives are in $S$. This means $S$ contains at most $2+(n-2)=n$ elements.

 Thus we have proved that $|S|\leq n=\lceil n/2\rceil$ for even $n$, and we are done.

 



\textbf{Solution 2}

Let $Z_n$ be the set of subsets satisfying the $a+b+c \neq 0$ condition for $n$, and let $s_n$ be the largest size of a set in $Z_n$. Let $n = 2k$ if $n$ is even, and $n = 2k-1$ if $n$ is odd. We note that $s_n \ge 2k$ due to the following constuction: 
 \[\{-2k+1, -2k + 3, \ldots, 2k-3, 2k-1\}\] or all of the odd numbers in the set. Then the sum of any three will be odd and thus nonzero. 

 
\textbf{Lemma 1}: $s_{n+1} \le s_{n} + 2$. If $M_{n+1} \in Z_{n+1}, |M_{n+1}| = s_{n+1}$, then we note that $M_{n+1} \setminus \{-2n+1, 2n-1\} \in Z_n$, so $|M_{n+1}| - 2 \le s_n$. $\square$

 
\textbf{Lemma 2}: $s_{2k} = s_{2k-1}$. Suppose, for sake of contradiction, that $M_{2k} \in Z_{n+1}$ and $M_{2k} > s_{2k-1}$. Remove $2k,-2k$ from $M_{2k}$, and partition the rest of the elements into two sets $P_{2k-1}, Q_{2k-1}$, where $P$ and $Q$ contain all of the positive and negative elements of $M_{2k}$, respectively. (obviously $0 \not \in M_{2k}$, because $0+0+0 = 0$). WLOG, suppose $|P_{2k-1}| \ge |Q_{2k-1}|$. Then $|P_{2k-1}| + |Q_{2k-1}| \ge s_{2k-1} - 1 \ge 2k - 1$. We now show the following two sub-results:

 


 
            \begin{enumerate}[label=(\roman*)]
                \item \textit{Sub-lemma (A)}: if $|P_{2k-1}| \ge k$, $-2k \not \in M_{2k}$ [and similar for $Q$]; and
            \end{enumerate}
            


 
            \begin{enumerate}[label=(\roman*)]
                \item \textit{Sub-lemma (B)}: we cannot have both $|P_{2k-1}| + |Q_{2k-1}| \ge 2k$ and $|Q_{2k-1}| < k$ simultaneously hold.
            \end{enumerate}
            
This is sufficient, because the only two elements that may be in $M_{2k}$ that are not in $P_{2k-1}, Q_{2k-1}$ are $2k$ and $-2k$; for $M_{2k} > s_{2k-1}$, we must either have $|P_{2k-1}| + |Q_{2k-1}| = s_{2k-1} - 1 \ge 2k-1$ and both $2k,-2k \in M_{2k}$ [but by pigeonhole $|P_{2k-1}| \ge k$, see sub-lemma (A)], or $|P_{2k-1}| + |Q_{2k-1}| = s_{2k-1} \ge 2k$, and $2k \in M_{2k}$, in which case by (A) we must have $|Q_{2k-1}| < k$, violating (B).

 
(A): Partition $\{1,2,\ldots,2k-1\}$ into the $k$ sets $\{k\}, \{1,2k-1\}, \{2, 2k-2\}, \cdots, \{k-1, k+1\}$. Because $(k+\delta) + (k-\delta) - 2k = 0$, then if any of those sets are within $P_{2k-1}$, $-2k \not \in M_{2k}$. But by Pigeonhole at most $k-1$ elements may be in $P_{2k-1}$, contradiction. 

 
(B): We prove this statement with another induction. We see that the statement easily holds true for $k=1$ or $2$, so suppose it is true for $k-1$, but [for sake of contradiction] false for $k$. Let $P_{2k-3} = P_{2k-1} \setminus \{2k-2, 2k-1\}$, and similarly for $Q_{2k-3}$. Again WLOG $|P_{2k-3}| \ge |Q_{2k-3}|$. Then we have $|P_{2k-1}| + |Q_{2k-1}| \le |P_{2k-3}| + |Q_{2k-3}| + 4$.

 
                \begin{itemize}
                    \item If $|P_{2k-3}| + |Q_{2k-3}| \ge 2k-2$, then by inductive hypothesis, we must have $|P_{2k-3}| = |Q_{2k-3}| = k-1$. But (A) implies that we cannot add $2k-2$ or $-2k+2$. So to satisfy $|P_{2k-1}| = |Q_{2k-1}| \ge 2k$ we must have $2k-1, -2k+1$ added, but then $|P_{2k-1}| = |Q_{2k-1}| = k$ contradiction.
                \end{itemize}


                
                \begin{itemize}
                    \item If $|P_{2k-3}| + |Q_{2k-3}| = 2k-3$, then at least three of $\{2k-2, 2k-1, -2k+2, -2k + 1\}$ added. But $|P_{2k-3}| \ge k-1$, and by (A) we have that $-2k+2$ cannot be added. If $|P_{2k-3}| \ge k$, then another grouping similar to (A) shows that $-2k+1$ canno be added, contradiction. So $|P_{2k-3}| = k-1$, $|Q_{2k-3}| = k-2$, and adding the three remaining elements gives $|P_{2k-1}| = |Q_{2k-1}| = k$ contradiction.
                \end{itemize}


                
                \begin{itemize}
                    \item If $|P_{2k-3}| + |Q_{2k-3}| = 2k-4$, then all four of $\{2k-2, 2k-1, -2k+2, -2k + 1\}$ must be added, and furthermore $|P_{2k-3}| > |Q_{2k-3}|$. Then $|P_{2k-3}| \ge k-1$, and by previous paragraph we cannot add $-2k+2$. $\square$
                \end{itemize}


                
So we have $\begin{cases} s_1 &= 2 \\ s_{2k} &= s_{2k-1} \\ s_{2k-1} &\le s_{2k-2} + 2 \\  \end{cases}$ and by induction, that $s_n \le \boxed{2\left\lceil \frac n2 \right \rceil} = 2k$, which we showed is achievable above.

 


	\item (\textit{USAMO 2009}) We define a \textit{chessboard polygon} to be a polygon whose sides are situated along lines of the form $x = a$ or $y = b$, where $a$ and $b$ are integers. These lines divide the interior into unit squares, which are shaded alternately grey and white so that adjacent squares have different colors. To tile a chessboard polygon by dominoes is to exactly cover the polygon by non-overlapping $1 \times 2$ rectangles. Finally, a \textit{tasteful tiling} is one which avoids the two configurations of dominoes shown on the left below. Two tilings of a $3 \times 4$ rectangle are shown; the first one is tasteful, while the second is not, due to the vertical dominoes in the upper right corner.


 
                    \begin{figure}[H]
                        \centering
                        \includegraphics[width=0.5\textwidth]{Problem/USAMO_I_2009_3_Problem_Figure_1.png}
                    \end{figure}
                    

  a) Prove that if a chessboard polygon can be tiled by dominoes, then it can be done so tastefully.


 b) Prove that such a tasteful tiling is unique.


 



\textbf{Solution}

<i>This problem needs a solution. If you have a solution for it, please help us out by \href{https://artofproblemsolving.com/wiki/index.php?title=2009\_USAMO\_Problems/Problem\_3&action=edit}{adding it}.</i> 

 


	\item (\textit{USAMO 2009}) For $n \ge 2$ let $a_1$, $a_2$, ..., $a_n$ be positive real numbers such that 


 $(a_1+a_2+ ... +a_n)\left( {1 \over a_1} + {1 \over a_2} + ... +{1 \over a_n} \right) \le \left(n+ {1 \over 2} \right) ^2$

 Prove that $\text{max}\, (a_1, a_2, ... ,a_n) \le  4\, \text{min}\, (a_1, a_2, ... , a_n)$.


 



\textbf{Solution 1}

Assume without loss of generality that $a_1 \geq a_2 \geq \cdots \geq a_n$. Now we seek to prove that $a_1 \le 4a_n$.

 By the \href{https://artofproblemsolving.com/wiki/index.php?title=Cauchy-Schwarz\_Inequality}{Cauchy-Schwarz Inequality}, 
 \begin{align*} (a_n+a_2+ a_3 + ... +a_{n-1}+a_1)\left({1 \over a_1} + {1 \over a_2} + ... +{1 \over a_n}\right) &\ge \left( \sqrt{a_n \over a_1} + n-2 + \sqrt{a_1 \over a_n} \right)^2 \\ (n+ {1 \over 2})^2 &\ge \left( \sqrt{a_n \over a_1} + n-2 + \sqrt{a_1 \over a_n} \right)^2 \\   n+ {1 \over 2} &\ge n-2 + \sqrt{a_n \over a_1} + \sqrt{a_1 \over a_n} \\ {5 \over 2} &\ge \sqrt{a_n \over a_1} + \sqrt{a_1 \over a_n} \\ {17 \over 4} &\ge {a_n \over a_1} + {a_1 \over a_n} \\  0 &\ge (a_1 - 4a_n)\left(a_1 - {a_n \over 4}\right) \end{align*}
Since $a_1 \ge a_n$, clearly $(a_1 - {a_n \over 4}) > 0$, dividing yields:

 
 \[0 \ge (a_1 - 4a_n) \Longrightarrow 4a_n \ge a_1\] 

 as desired.

 



\textbf{Solution 2}

Assume without loss of generality that $a_1 \geq a_2 \geq \cdots \geq a_n$.Using the Cauchy–Bunyakovsky–Schwarz inequality and the inequality given, 
 \begin{align*} &\left(a_1+a_2 + ... +a_n + 3a_n -\frac{3a_1}{4}\right)\left({1 \over a_1} + {1 \over a_2} + ... +{1 \over a_n}\right) \\ =&\left(\frac{a_1}{4}+a_2 + ... +a_{n-1}+4a_n\right)\left({1 \over a_1} + {1 \over a_2} + ... +{1 \over a_n}\right) \\ \ge& \left(\frac{1}{2}+n-2+2\right)^2 \\ =&\left(n+\frac{1}{2}\right)^2 \\ \ge& (a_1+a_2 + ... +a_{n})\left({1 \over a_1} + {1 \over a_2} + ... +{1 \over a_n}\right).\end{align*}
(Note that $n-2 \ge 0$ since $n \ge 2$ as given!)This implies that $3a_n -\frac{3a_1}{4} \ge 0 \iff 4a_n \ge a_1$ as desired.

 


	\item (\textit{USAMO 2009}) Trapezoid $ABCD$, with $\overline{AB}||\overline{CD}$, is inscribed in circle $\omega$ and point $G$ lies inside triangle $BCD$.  Rays $AG$ and $BG$ meet $\omega$ again at points $P$ and $Q$, respectively.  Let the line through $G$ parallel to $\overline{AB}$ intersect $\overline{BD}$ and $\overline{BC}$ at points $R$ and $S$, respectively.  Prove that quadrilateral $PQRS$ is cyclic if and only if $\overline{BG}$ bisects $\angle CBD$.


 



\textbf{Solution 1}

We will use directed angles in this solution. Extend $QR$ to $T$ as follows:

 
                    \begin{figure}[H]
                        \centering
                        \includegraphics[width=0.5\textwidth]{Solution/USAMO_I_2009_5_Solution_Figure_1.png}
                    \end{figure}
                    

 \textbf{If}:

 Note that 
 \begin{align*}\measuredangle GBT+\measuredangle TRG&=\frac{m\widehat{TQ}}{2}+\measuredangle TRB+\measuredangle BRG\\ &=\frac{m\widehat{TQ}+m\widehat{DQ}+m\widehat{CB}+m\widehat{BT}}{2}.\\ \end{align*}
Thus, $BTRG$ is cyclic.

 Also, note that $GSCP$ is cyclic because 
 \begin{align*}\measuredangle CSG+\measuredangle GPC&=\measuredangle CBA+\measuredangle APC\\ &=180^\circ\text{ or }0^\circ, \end{align*}
 depending on the configuration.

 Next, we have $T, G, C$ are collinear since 
 \[\measuredangle GTR=\measuredangle GBR=\frac{m\widehat{DQ}}{2}=\frac{m\widehat{QC}}{2}=\measuredangle CTQ.\]

 Therefore, 
 \begin{align*}\measuredangle RQP+\measuredangle PSR&=\frac{m\widehat{PBT}}{2}+\measuredangle PCG\\ &=\frac{m\widehat{PBT}+m\widehat{TDP}}{2}\\ &=180^\circ \end{align*}
, so $PQRS$ is cyclic.

 \textbf{Only If}:These steps can be reversed.

 



\textbf{Solution 2 (Projective)}

Extend $QR$ to $T$, and let line $l \parallel AB$ intersect $\omega$ at $K$ and another point $V$, as shown:

 
                    \begin{figure}[H]
                        \centering
                        \includegraphics[width=0.5\textwidth]{Solution/USAMO_I_2009_5_Solution_Figure_2.png}
                    \end{figure}
                    

 \textbf{If}:

 Suppose that $VP \cap CB = S'$, and $AC \cap QV = R'$. Pascal's theorem on the tuple $(V, P, A, C, B, Q)$ implies that the points $S'$, $R'$, and $G = PA \cap BQ$ are collinear. However, $AC$ and $BD$ are symmetrical with respect to the axis of symmetry of trapezoid $ABCD$, and $TQ$ and $VQ$ are also symmetrical with respect to the axis of symmetry of $ABCD$ (as $Q$ is the midpoint of $\overset{\frown}{DC}$, and $TV \parallel DC$). Since $R = BD \cap TQ$, $R$ and $R'$ are symmetric with respect to the axis of symmetry of trapezoid $ABCD$. This implies that line $R'G$ is equivalent to line $RG$. Thus, $S'$ lies on line $RG$. However, $S = BC \cap RG$, so this implies that $S' = S$.

 Now note that $TVPQ$ is cyclic. Since $TV \parallel RS$, $\measuredangle VTQ = \measuredangle SRQ$. However, $\measuredangle VTQ + \measuredangle VPQ = 180^{\circ} = \measuredangle SRQ + \measuredangle SPQ$. Therefore, $PQRS$ is cyclic.

 \textbf{Only If}:

 Consider the same setup, except $Q$ is no longer the midpoint of $\overset{\frown}{DC}$. Note that $TV$ must be parallel to $RG$ in order for $PQRS$  to be cyclic. We claim that $S' = S$ and hope to reach a contradiction. Pascal's theorem on the tuple $(V, P, A, C, B, Q)$ implies that $S'$, $R'$, and $G = PA \cap BQ$ are collinear. However, there exists a unique point $Q$ such that $HQ$, $AC$, and $RG$ are concurrent. By \textbf{If}, $Q$ must be the midpoint of $\overset{\frown}{DC}$ in order for the concurrency to occur; hence, $R' \notin RS$. Then $R'G \cap BC = S' \neq S$, since $RG \cap BC = S$. However, this is a contradiction, so therefore $TV$ cannot be parallel to $RG$ and $PQRS$ is not cyclic.

 Solution by \textbf{TheBoomBox77}

 


	\item (\textit{USAMO 2009}) Let $s_1, s_2, s_3, \ldots$ be an infinite, nonconstant sequence of rational numbers, meaning it is not the case that $s_1 = s_2 = s_3 = \cdots.$  Suppose that $t_1, t_2, t_3, \ldots$ is also an infinite, nonconstant sequence of rational numbers with the property that $(s_i - s_j)(t_i - t_j)$ is an integer for all $i$ and $j$.  Prove that there exists a rational number $r$ such that $(s_i - s_j)r$ and $(t_i - t_j)/r$ are integers for all $i$ and $j$.


 



\textbf{Solution}

Suppose the $s_i$ can be represented as $\frac{a_i}{b_i}$ for every $i$, and suppose $t_i$ can be represented as $\frac{c_i}{d_i}$. Let's start with only the first two terms in the two sequences, $s_1$ and $s_2$ for sequence $s$ and $t_1$ and $t_2$ for sequence $t$. Then by the conditions of the problem, we have $(s_2 - s_1)(t_2 - t_1)$ is an integer, or $(\frac{a_2}{b_2} - \frac{a_1}{b_1})(\frac{c_2}{d_2} - \frac{c_1}{d_1})$ is an integer. Now we can set $r = \frac{b_1 b_2}{d_1 d_2}$, because the least common denominator of $s_2 - s_1$ is $b_1 b_2$ and of $t_2 - t_1$ is $d_1 d_2$, and multiplying or dividing appropriately by $\frac{b_1 b_2}{d_1 d_2}$ will always give an integer. 

 Now suppose we kept adding $s_i$ and $t_i$ until we get to $s_m = \frac{a_m}{b_m}$ in sequence $s$ and $t_m = \frac{c_m}{d_m}$ in sequence $t$ so that $(t_m - t_i)(s_m - s_i)$ is an integer for all $i$ with $1 \le i < m$, where $m$ is a positive integer. At this point, we will have $r$ = $\frac{\prod_{n=1}^{m}b_n}{\prod_{n=1}^{m}d_n}$, because these are the least common denominators of the two sequences up to $m$. As we keep adding $s_i$ and $t_i$, $r$ will always have value $\frac{\prod_{n=1}^{m}b_n}{\prod_{n=1}^{m}d_n}$, and we are done.

 


\end{enumerate}